% Generated by: cargo run -p jankurai -- paper public-repo-scores
% Source: paper/data/public-repo-scores-20260506T014156Z.json
% Command: cargo run -p jankurai -- paper public-repo-scores --source paper/data/public-repo-scores-20260506T014156Z.json --out paper/tex/generated/public_repo_score_tables.tex
% DO NOT EDIT BY HAND.
% Source artifact: paper/data/public-repo-scores-20260506T014156Z.json
% Generated at: 2026-05-06T01:41:56Z
% Jankurai version: 0.8.8

\newcommand{\PublicRepoScoreTopTable}{%
\begin{table*}[t]
\caption{Star-rank ordered advisory public-repository sample. Scores are advisory posture signals, not certification results.}
\label{tab:public-repo-top-scores}
\centering
\scriptsize
\setlength{\tabcolsep}{3pt}
\begin{tabularx}{\textwidth}{R{0.035\textwidth} L{0.205\textwidth} R{0.060\textwidth} R{0.050\textwidth} R{0.065\textwidth} R{0.070\textwidth} L{0.150\textwidth} Y}
\toprule
\JKTableHeader
\textbf{Rank} & \textbf{Repo} & \textbf{Stars} & \textbf{Score} & \textbf{Issues} & \textbf{Hard/Soft} & \textbf{Severity} & \textbf{Top categories} \\
\midrule
\JKDenseRows
1 & \path{rustdesk/rustdesk} & 114k & \JKScoreCell{29} & 1,345 & 1,333/12 & med 12; high 1,332; crit 1 & security 684; vibe 617; python 17 \\
2 & \path{denoland/deno} & 107k & \JKScoreCell{42} & 3,130 & 3,116/14 & med 14; high 3,115; crit 1 & vibe 1,588; security 1,455; boundary 52 \\
3 & \path{tauri-apps/tauri} & 106k & \JKScoreCell{38} & 485 & 473/12 & med 12; high 472; crit 1 & vibe 280; security 141; boundary 43 \\
4 & \path{astral-sh/uv} & 84k+ & \JKScoreCell{38} & 1,258 & 1,246/12 & med 12; high 1,245; crit 1 & vibe 1,111; security 101; python 14 \\
5 & \path{zed-industries/zed} & 81.7k & \JKScoreCellBest{47} & 3,328 & 3,315/13 & med 13; high 3,314; crit 1 & vibe 2,004; security 1,225; data 68 \\
6 & \path{alacritty/alacritty} & 63.9k & \JKScoreCell{33} & 231 & 219/12 & med 12; high 219 & security 127; vibe 89; context 3 \\
7 & \path{BurntSushi/ripgrep} & 63.3k & \JKScoreCell{30} & 80 & 68/12 & med 12; high 68 & vibe 47; security 13; proof 3 \\
8 & \path{meilisearch/meilisearch} & 57.4k & \JKScoreCell{41} & 793 & 782/11 & med 11; high 781; crit 1 & vibe 671; security 102; proof 3 \\
9 & \path{typst/typst} & 53.3k & \JKScoreCell{32} & 398 & 386/12 & med 12; high 386 & vibe 358; security 21; proof 3 \\
10 & \path{astral-sh/ruff} & 47.4k & \JKScoreCell{44} & 2,842 & 2,829/13 & med 13; high 2,828; crit 1 & vibe 2,724; python 35; security 30 \\
\bottomrule
\end{tabularx}
\JKResetRows
\end{table*}
}

\newcommand{\PublicRepoScoreAggregateTable}{%
\begin{table}[t]
\caption{Aggregate advisory scan posture.}
\label{tab:public-repo-aggregate}
\centering
\scriptsize
\setlength{\tabcolsep}{4pt}
\begin{tabularx}{\columnwidth}{L{0.60\columnwidth} R{0.28\columnwidth}}
\toprule
\JKTableHeader
\textbf{Metric} & \textbf{Value} \\
\midrule
\JKDenseRows
Repositories scanned & 30 \\
Successful scans & 30 \\
Failed scans & 0 \\
Minimum score & 14 \\
Maximum score & 47 (best observed) \\
Average score & 33.0 \\
Median score (upper middle) & 33 \\
Findings total & 20,024 \\
Hard findings total & 19,651 \\
Hard findings share & 98.1\% \\
\bottomrule
\end{tabularx}
\JKResetRows
\end{table}
}

\newcommand{\PublicRepoWeakDimensionTable}{%
\begin{table}[t]
\caption{Most frequent top finding categories in the advisory scan.}
\label{tab:public-repo-weak-dimensions}
\centering
\scriptsize
\setlength{\tabcolsep}{4pt}
\begin{tabularx}{\columnwidth}{Y R{0.18\columnwidth}}
\toprule
\JKTableHeader
\textbf{Category} & \textbf{Findings} \\
\midrule
\JKDenseRows
vibe & 14,628 \\
security & 4,232 \\
boundary & 207 \\
python & 107 \\
data & 68 \\
proof & 51 \\
\bottomrule
\end{tabularx}
\JKResetRows
\end{table}
}

\newcommand{\PublicRepoScoreAppendixTable}{%
\begingroup
\scriptsize
\setlength{\tabcolsep}{3pt}
\setlength{\LTleft}{0pt}
\setlength{\LTright}{0pt}
\JKDenseRows
\begin{longtable}{@{}R{0.035\textwidth} L{0.180\textwidth} R{0.055\textwidth} R{0.045\textwidth} R{0.060\textwidth} R{0.070\textwidth} L{0.145\textwidth} L{0.300\textwidth}@{}}
\caption{Full 30-repository advisory scoring run.}\label{tab:public-repo-full}\\
\toprule
\JKTableHeader
\textbf{Rank} & \textbf{Repo} & \textbf{Stars} & \textbf{Score} & \textbf{Issues} & \textbf{Hard/Soft} & \textbf{Severity} & \textbf{Top categories} \\
\midrule
\endfirsthead
\toprule
\JKTableHeader
\textbf{Rank} & \textbf{Repo} & \textbf{Stars} & \textbf{Score} & \textbf{Issues} & \textbf{Hard/Soft} & \textbf{Severity} & \textbf{Top categories} \\
\midrule
\endhead
1 & \path{rustdesk/rustdesk} & 114k & \JKScoreCell{29} & 1,345 & 1,333/12 & med 12; high 1,332; crit 1 & security 684; vibe 617; python 17 \\
2 & \path{denoland/deno} & 107k & \JKScoreCell{42} & 3,130 & 3,116/14 & med 14; high 3,115; crit 1 & vibe 1,588; security 1,455; boundary 52 \\
3 & \path{tauri-apps/tauri} & 106k & \JKScoreCell{38} & 485 & 473/12 & med 12; high 472; crit 1 & vibe 280; security 141; boundary 43 \\
4 & \path{astral-sh/uv} & 84k+ & \JKScoreCell{38} & 1,258 & 1,246/12 & med 12; high 1,245; crit 1 & vibe 1,111; security 101; python 14 \\
5 & \path{zed-industries/zed} & 81.7k & \JKScoreCellBest{47} & 3,328 & 3,315/13 & med 13; high 3,314; crit 1 & vibe 2,004; security 1,225; data 68 \\
6 & \path{alacritty/alacritty} & 63.9k & \JKScoreCell{33} & 231 & 219/12 & med 12; high 219 & security 127; vibe 89; context 3 \\
7 & \path{BurntSushi/ripgrep} & 63.3k & \JKScoreCell{30} & 80 & 68/12 & med 12; high 68 & vibe 47; security 13; proof 3 \\
8 & \path{meilisearch/meilisearch} & 57.4k & \JKScoreCell{41} & 793 & 782/11 & med 11; high 781; crit 1 & vibe 671; security 102; proof 3 \\
9 & \path{typst/typst} & 53.3k & \JKScoreCell{32} & 398 & 386/12 & med 12; high 386 & vibe 358; security 21; proof 3 \\
10 & \path{astral-sh/ruff} & 47.4k & \JKScoreCell{44} & 2,842 & 2,829/13 & med 13; high 2,828; crit 1 & vibe 2,724; python 35; security 30 \\
11 & \path{rtk-ai/rtk} & 41.8k & \JKScoreCell{25} & 451 & 437/14 & med 14; high 436; crit 1 & vibe 371; agent 41; security 11 \\
12 & \path{vercel-labs/agent-browser} & 31.7k & \JKScoreCell{38} & 226 & 215/11 & med 11; high 214; crit 1 & vibe 165; security 34; boundary 8 \\
13 & \path{zeroclaw-labs/zeroclaw} & 31.0k & \JKScoreCell{40} & 1,317 & 1,304/13 & med 13; high 1,303; crit 1 & vibe 1,192; security 70; boundary 23 \\
14 & \path{googleworkspace/cli} & 25.8k & \JKScoreCell{40} & 114 & 103/11 & med 11; high 102; crit 1 & vibe 75; security 19; proof 3 \\
15 & \path{AlexsJones/llmfit} & 25.3k & \JKScoreCell{35} & 110 & 97/13 & med 13; high 96; crit 1 & vibe 80; security 9; proof 3 \\
16 & \path{RightNow-AI/openfang} & 17.2k & \JKScoreCell{42} & 695 & 682/13 & med 13; high 681; crit 1 & vibe 665; security 9; proof 4 \\
17 & \path{xai-org/x-algorithm} & 16.4k & \JKScoreCell{24} & 38 & 25/13 & med 13; high 25 & vibe 12; proof 4; security 3 \\
18 & \path{microsoft/RustTraining} & 14.2k & \JKScoreCell{31} & 34 & 21/13 & med 13; high 20; crit 1 & vibe 9; proof 5; security 5 \\
19 & \path{nearai/ironclaw} & 12.1k & \JKScoreCell{42} & 1,968 & 1,955/13 & med 13; high 1,954; crit 1 & vibe 1,768; security 119; python 30 \\
20 & \path{h4ckf0r0day/obscura} & 10.3k & \JKScoreCell{25} & 63 & 51/12 & med 12; high 51 & vibe 35; security 8; proof 4 \\
21 & \path{octasoft-ltd/wsl-ui} & 178 & \JKScoreCell{30} & 321 & 309/12 & med 12; high 308; crit 1 & vibe 288; boundary 8; security 6 \\
22 & \path{kvnxiao/tauri-tanstack-start-react-template} & 37 & \JKScoreCell{39} & 24 & 12/12 & med 12; high 12 & proof 3; security 3; context 3 \\
23 & \path{fudanglp/tauri-fastapi-full-stack-template} & 17 & \JKScoreCell{23} & 87 & 74/13 & med 13; high 73; crit 1 & vibe 41; python 11; security 8 \\
24 & \path{sergioadevita/notemac-plus-plus} & 9 & \JKScoreCell{23} & 344 & 330/14 & med 14; high 329; crit 1 & vibe 264; boundary 51; security 7 \\
25 & \path{ianho7/maptoposter-online} & 7 & \JKScoreCell{14} & 66 & 53/13 & med 13; high 53 & vibe 32; boundary 8; proof 4 \\
26 & \path{exit-zero-labs/threat-forge} & 6 & \JKScoreCell{32} & 142 & 129/13 & med 13; high 128; crit 1 & vibe 107; boundary 10; security 6 \\
27 & \path{lostf1sh/rustune} & 5 & \JKScoreCell{23} & 38 & 26/12 & med 12; high 26 & vibe 12; security 5; proof 4 \\
28 & \path{Duri686/RustQuantLab} & 4 & \JKScoreCell{26} & 45 & 33/12 & med 12; high 32; crit 1 & vibe 17; security 5; boundary 4 \\
29 & \path{kevinpiao1025/tauri-react-typescript-tailwind} & 3 & \JKScoreCell{39} & 21 & 10/11 & med 11; high 10 & proof 4; security 3; context 3 \\
30 & \path{MarkShawn2020/lovtauri} & 2 & \JKScoreCell{26} & 30 & 18/12 & med 12; high 18 & vibe 6; proof 4; security 3 \\
\bottomrule
\end{longtable}
\JKResetRows
\endgroup
}

\newcommand{\PublicRepoScoreAppendixPlot}{%
\begin{figure*}[t]
\centering
\definecolor{JKPlotRed}{HTML}{B54B4B}
\definecolor{JKPlotAmber}{HTML}{D8A33A}
\definecolor{JKPlotGreen}{HTML}{2F8F5B}
\resizebox{\textwidth}{!}{%
\begin{tikzpicture}[x=0.36cm,y=0.050cm]
\fill[JKHeader!55] (-0.85,-31) rectangle (25.25,106);
\draw[JKLine] (-0.45,0) -- (24.95,0);
\draw[JKLine] (-0.45,50) -- (24.95,50);
\draw[JKLine] (-0.45,85) -- (24.95,85);
\draw[JKLine] (-0.45,100) -- (24.95,100);
\node[anchor=east,scale=0.72,text=JKNavy] at (-0.55,0) {0};
\node[anchor=east,scale=0.72,text=JKNavy] at (-0.55,50) {50};
\node[anchor=east,scale=0.72,text=JKNavy] at (-0.55,85) {85};
\node[anchor=east,scale=0.72,text=JKNavy] at (-0.55,100) {100};
\node[anchor=south,scale=0.78,text=JKNavy] at (12.2,103) {Advisory Jankurai score};
\filldraw[draw=JKNavy!45,fill=JKPlotGreen] (0.00,0) rectangle (0.52,100);
\node[anchor=south,rotate=90,scale=0.62,text=JKNavy] at (0.26,101) {100};
\node[anchor=east,rotate=45,scale=0.52,text=black] at (0.40,-3) {\textbf{Jankurai}};
\filldraw[draw=JKNavy!45,fill=JKPlotRed!94!JKPlotAmber] (0.78,0) rectangle (1.30,47);
\node[anchor=south,rotate=90,scale=0.62,text=JKNavy] at (1.04,48) {47};
\node[anchor=east,rotate=45,scale=0.52,text=black] at (1.18,-3) {\texttt{zed-industries/zed}};
\filldraw[draw=JKNavy!45,fill=JKPlotRed!88!JKPlotAmber] (1.56,0) rectangle (2.08,44);
\node[anchor=south,rotate=90,scale=0.62,text=JKNavy] at (1.82,45) {44};
\node[anchor=east,rotate=45,scale=0.52,text=black] at (1.96,-3) {\texttt{astral-sh/ruff}};
\filldraw[draw=JKNavy!45,fill=JKPlotRed!84!JKPlotAmber] (2.34,0) rectangle (2.86,42);
\node[anchor=south,rotate=90,scale=0.62,text=JKNavy] at (2.60,43) {42};
\node[anchor=east,rotate=45,scale=0.52,text=black] at (2.74,-3) {\texttt{denoland/deno}};
\filldraw[draw=JKNavy!45,fill=JKPlotRed!84!JKPlotAmber] (3.12,0) rectangle (3.64,42);
\node[anchor=south,rotate=90,scale=0.62,text=JKNavy] at (3.38,43) {42};
\node[anchor=east,rotate=45,scale=0.52,text=black] at (3.52,-3) {\texttt{nearai/ironclaw}};
\filldraw[draw=JKNavy!45,fill=JKPlotRed!84!JKPlotAmber] (3.90,0) rectangle (4.42,42);
\node[anchor=south,rotate=90,scale=0.62,text=JKNavy] at (4.16,43) {42};
\node[anchor=east,rotate=45,scale=0.52,text=black] at (4.30,-3) {\texttt{RightNow-AI/openfang}};
\filldraw[draw=JKNavy!45,fill=JKPlotRed!82!JKPlotAmber] (4.68,0) rectangle (5.20,41);
\node[anchor=south,rotate=90,scale=0.62,text=JKNavy] at (4.94,42) {41};
\node[anchor=east,rotate=45,scale=0.52,text=black] at (5.08,-3) {\texttt{meilisearch/meilisearch}};
\filldraw[draw=JKNavy!45,fill=JKPlotRed!80!JKPlotAmber] (5.46,0) rectangle (5.98,40);
\node[anchor=south,rotate=90,scale=0.62,text=JKNavy] at (5.72,41) {40};
\node[anchor=east,rotate=45,scale=0.52,text=black] at (5.86,-3) {\texttt{googleworkspace/cli}};
\filldraw[draw=JKNavy!45,fill=JKPlotRed!80!JKPlotAmber] (6.24,0) rectangle (6.76,40);
\node[anchor=south,rotate=90,scale=0.62,text=JKNavy] at (6.50,41) {40};
\node[anchor=east,rotate=45,scale=0.52,text=black] at (6.64,-3) {\texttt{zeroclaw-labs/zeroclaw}};
\filldraw[draw=JKNavy!45,fill=JKPlotRed!78!JKPlotAmber] (7.02,0) rectangle (7.54,39);
\node[anchor=south,rotate=90,scale=0.62,text=JKNavy] at (7.28,40) {39};
\node[anchor=east,rotate=45,scale=0.52,text=black] at (7.42,-3) {\texttt{kevinpiao1025/tauri-react-typescript-tailwind}};
\filldraw[draw=JKNavy!45,fill=JKPlotRed!78!JKPlotAmber] (7.80,0) rectangle (8.32,39);
\node[anchor=south,rotate=90,scale=0.62,text=JKNavy] at (8.06,40) {39};
\node[anchor=east,rotate=45,scale=0.52,text=black] at (8.20,-3) {\texttt{kvnxiao/tauri-tanstack-start-react-template}};
\filldraw[draw=JKNavy!45,fill=JKPlotRed!76!JKPlotAmber] (8.58,0) rectangle (9.10,38);
\node[anchor=south,rotate=90,scale=0.62,text=JKNavy] at (8.84,39) {38};
\node[anchor=east,rotate=45,scale=0.52,text=black] at (8.98,-3) {\texttt{astral-sh/uv}};
\filldraw[draw=JKNavy!45,fill=JKPlotRed!76!JKPlotAmber] (9.36,0) rectangle (9.88,38);
\node[anchor=south,rotate=90,scale=0.62,text=JKNavy] at (9.62,39) {38};
\node[anchor=east,rotate=45,scale=0.52,text=black] at (9.76,-3) {\texttt{tauri-apps/tauri}};
\filldraw[draw=JKNavy!45,fill=JKPlotRed!76!JKPlotAmber] (10.14,0) rectangle (10.66,38);
\node[anchor=south,rotate=90,scale=0.62,text=JKNavy] at (10.40,39) {38};
\node[anchor=east,rotate=45,scale=0.52,text=black] at (10.54,-3) {\texttt{vercel-labs/agent-browser}};
\filldraw[draw=JKNavy!45,fill=JKPlotRed!70!JKPlotAmber] (10.92,0) rectangle (11.44,35);
\node[anchor=south,rotate=90,scale=0.62,text=JKNavy] at (11.18,36) {35};
\node[anchor=east,rotate=45,scale=0.52,text=black] at (11.32,-3) {\texttt{AlexsJones/llmfit}};
\filldraw[draw=JKNavy!45,fill=JKPlotRed!66!JKPlotAmber] (11.70,0) rectangle (12.22,33);
\node[anchor=south,rotate=90,scale=0.62,text=JKNavy] at (11.96,34) {33};
\node[anchor=east,rotate=45,scale=0.52,text=black] at (12.10,-3) {\texttt{alacritty/alacritty}};
\filldraw[draw=JKNavy!45,fill=JKPlotRed!64!JKPlotAmber] (12.48,0) rectangle (13.00,32);
\node[anchor=south,rotate=90,scale=0.62,text=JKNavy] at (12.74,33) {32};
\node[anchor=east,rotate=45,scale=0.52,text=black] at (12.88,-3) {\texttt{exit-zero-labs/threat-forge}};
\filldraw[draw=JKNavy!45,fill=JKPlotRed!64!JKPlotAmber] (13.26,0) rectangle (13.78,32);
\node[anchor=south,rotate=90,scale=0.62,text=JKNavy] at (13.52,33) {32};
\node[anchor=east,rotate=45,scale=0.52,text=black] at (13.66,-3) {\texttt{typst/typst}};
\filldraw[draw=JKNavy!45,fill=JKPlotRed!62!JKPlotAmber] (14.04,0) rectangle (14.56,31);
\node[anchor=south,rotate=90,scale=0.62,text=JKNavy] at (14.30,32) {31};
\node[anchor=east,rotate=45,scale=0.52,text=black] at (14.44,-3) {\texttt{microsoft/RustTraining}};
\filldraw[draw=JKNavy!45,fill=JKPlotRed!60!JKPlotAmber] (14.82,0) rectangle (15.34,30);
\node[anchor=south,rotate=90,scale=0.62,text=JKNavy] at (15.08,31) {30};
\node[anchor=east,rotate=45,scale=0.52,text=black] at (15.22,-3) {\texttt{BurntSushi/ripgrep}};
\filldraw[draw=JKNavy!45,fill=JKPlotRed!60!JKPlotAmber] (15.60,0) rectangle (16.12,30);
\node[anchor=south,rotate=90,scale=0.62,text=JKNavy] at (15.86,31) {30};
\node[anchor=east,rotate=45,scale=0.52,text=black] at (16.00,-3) {\texttt{octasoft-ltd/wsl-ui}};
\filldraw[draw=JKNavy!45,fill=JKPlotRed!58!JKPlotAmber] (16.38,0) rectangle (16.90,29);
\node[anchor=south,rotate=90,scale=0.62,text=JKNavy] at (16.64,30) {29};
\node[anchor=east,rotate=45,scale=0.52,text=black] at (16.78,-3) {\texttt{rustdesk/rustdesk}};
\filldraw[draw=JKNavy!45,fill=JKPlotRed!52!JKPlotAmber] (17.16,0) rectangle (17.68,26);
\node[anchor=south,rotate=90,scale=0.62,text=JKNavy] at (17.42,27) {26};
\node[anchor=east,rotate=45,scale=0.52,text=black] at (17.56,-3) {\texttt{Duri686/RustQuantLab}};
\filldraw[draw=JKNavy!45,fill=JKPlotRed!52!JKPlotAmber] (17.94,0) rectangle (18.46,26);
\node[anchor=south,rotate=90,scale=0.62,text=JKNavy] at (18.20,27) {26};
\node[anchor=east,rotate=45,scale=0.52,text=black] at (18.34,-3) {\texttt{MarkShawn2020/lovtauri}};
\filldraw[draw=JKNavy!45,fill=JKPlotRed!50!JKPlotAmber] (18.72,0) rectangle (19.24,25);
\node[anchor=south,rotate=90,scale=0.62,text=JKNavy] at (18.98,26) {25};
\node[anchor=east,rotate=45,scale=0.52,text=black] at (19.12,-3) {\texttt{h4ckf0r0day/obscura}};
\filldraw[draw=JKNavy!45,fill=JKPlotRed!50!JKPlotAmber] (19.50,0) rectangle (20.02,25);
\node[anchor=south,rotate=90,scale=0.62,text=JKNavy] at (19.76,26) {25};
\node[anchor=east,rotate=45,scale=0.52,text=black] at (19.90,-3) {\texttt{rtk-ai/rtk}};
\filldraw[draw=JKNavy!45,fill=JKPlotRed!48!JKPlotAmber] (20.28,0) rectangle (20.80,24);
\node[anchor=south,rotate=90,scale=0.62,text=JKNavy] at (20.54,25) {24};
\node[anchor=east,rotate=45,scale=0.52,text=black] at (20.68,-3) {\texttt{xai-org/x-algorithm}};
\filldraw[draw=JKNavy!45,fill=JKPlotRed!46!JKPlotAmber] (21.06,0) rectangle (21.58,23);
\node[anchor=south,rotate=90,scale=0.62,text=JKNavy] at (21.32,24) {23};
\node[anchor=east,rotate=45,scale=0.52,text=black] at (21.46,-3) {\texttt{fudanglp/tauri-fastapi-full-stack-template}};
\filldraw[draw=JKNavy!45,fill=JKPlotRed!46!JKPlotAmber] (21.84,0) rectangle (22.36,23);
\node[anchor=south,rotate=90,scale=0.62,text=JKNavy] at (22.10,24) {23};
\node[anchor=east,rotate=45,scale=0.52,text=black] at (22.24,-3) {\texttt{lostf1sh/rustune}};
\filldraw[draw=JKNavy!45,fill=JKPlotRed!46!JKPlotAmber] (22.62,0) rectangle (23.14,23);
\node[anchor=south,rotate=90,scale=0.62,text=JKNavy] at (22.88,24) {23};
\node[anchor=east,rotate=45,scale=0.52,text=black] at (23.02,-3) {\texttt{sergioadevita/notemac-plus-plus}};
\filldraw[draw=JKNavy!45,fill=JKPlotRed!28!JKPlotAmber] (23.40,0) rectangle (23.92,14);
\node[anchor=south,rotate=90,scale=0.62,text=JKNavy] at (23.66,15) {14};
\node[anchor=east,rotate=45,scale=0.52,text=black] at (23.80,-3) {\texttt{ianho7/maptoposter-online}};
\end{tikzpicture}%
}
\caption{Advisory public-repository score ranking with the Jankurai repository baseline included. Bars are sorted by score; color moves from red through amber to green as scores approach the merge-witness-ready band.}
\label{fig:public-repo-score-ranking}
\end{figure*}
}


\section{Merge-Witness Readiness Scan}

The strongest field evidence in this edition is deliberately advisory. On May 6, 2026, Jankurai 0.8.8 scanned 30 public GitHub repositories against the current target profile. The paper uses a repo-local JSON source and digest receipt built from the run summary and star-rank table \cite{jankuraiPublicRepoRun2026}. This is not a certification result, not a defect claim, and not an empirical incident study. It is a posture scan asking whether prominent repositories expose merge-witness-ready ownership, proof routing, generated-zone receipts, context routing, proof freshness, and repair evidence.

The sample includes mature and widely used public projects such as Zed, Ruff, Deno, Meilisearch, Tauri, and ripgrep \cite{zedGithub2026,ruffGithub2026,denoGithub2026,meilisearchGithub2026,tauriGithub2026,ripgrepGithub2026}. Many of these projects have serious engineering cultures, active maintainers, and conventional CI. The surprising result is therefore not that small or immature repositories score low. It is that conventional maturity and popularity do not automatically produce the repository-alignment surfaces Jankurai expects.

The top observed public-repository score was 47, and every scanned public repository remained below the 50-point red band. The average score was 33.0. Across the run, Jankurai emitted 20,024 findings, including 19,651 hard findings, a hard-finding share of about 98.1\%. Those counts should be read as evidence gaps and repair opportunities under Jankurai policy, not as accusations against maintainers.

\PublicRepoScoreAggregateTable

The common missing surfaces were consistent with the paper's thesis. Repositories often had tests and CI, but did not expose an audit lane as a CI replacement artifact; did not bind owner maps, test maps, and generated-zone policy into a merge witness; did not leave fresh proof receipts for changed paths; and did not serialize repair queues in a form that a weaker implementation agent could consume. This matches the broader AI-assisted development evidence: DORA frames AI as an amplifier of existing delivery capability, not an automatic substitute for engineering discipline; security and secret-sprawl reports show plausible generated work can increase evidence risk; developer surveys and agent benchmarks show that nearly-correct output, environment setup, and context retrieval remain expensive failure modes \cite{doraStateAIAssistedSoftwareDevelopment2025,veracodeGenaiCodeSecurity2025,gitguardianSecretsSprawl2026,gitguardianSecretsSprawl2026Blog,stackOverflowAISurvey2025,sweAgentAci2024,setupBench2025,contextBench2026}.

\PublicRepoWeakDimensionTable

The category frequency makes the gap concrete. Vibe markers, security posture, boundary evidence, proof routing, Python containment, data access, and context routing dominate the top category summaries. These are not claims that the projects lack quality. They are claims that the repositories generally do not publish the specific merge-evidence interface required for a Jankurai witness.

The caveat matters. HLT marker counts, hard caps, and target-profile assumptions are strict policy proxies. A task marker, large file, Python placement, generated-boundary heuristic, or missing Jankurai CI artifact can be a false positive or a locally accepted tradeoff. The advisory scan is useful because it names repairable evidence gaps, not because it proves insecurity or nonconformance. Popularity, language choice, and conventional CI do not by themselves create repository alignment. The missing merge witness is exactly the standard gap Jankurai targets.

The full named table and a two-column score plot remain in Appendix~F. Keeping the main body aggregate-first avoids turning an exploratory readiness scan into a public quality leaderboard.
